\documentclass[11pt]{article}
\usepackage[margin=1in]{geometry}
\usepackage{amsmath,amssymb,mathtools}
\usepackage{booktabs,longtable,array}
\usepackage{enumitem}
\usepackage{xcolor}
\usepackage{listings}
\usepackage{tikz}
\usepackage{hyperref}
\usepackage{microtype}
\usetikzlibrary{arrows.meta,positioning,fit,shapes.geometric,calc}
\hypersetup{colorlinks=true,linkcolor=black,citecolor=black,urlcolor=blue}
\lstset{basicstyle=\ttfamily\small,frame=single,columns=fullflexible,keepspaces=true,showstringspaces=false,breaklines=true}
\newcommand{\R}{\mathbb{R}}
\newcommand{\Z}{\mathbb{Z}}
\newcommand{\N}{\mathbb{N}}
\newcommand{\Qdom}{\mathcal{Q}}
\newcommand{\Pin}{\mathcal{P}_{\mathrm{in}}}
\newcommand{\Yset}{\mathcal{Y}}
\newcommand{\Hull}{\operatorname{hull}}
\newcommand{\range}{\operatorname{range}}
\newcommand{\midp}{\operatorname{mid}}
\newcommand{\rad}{\operatorname{rad}}
\newcommand{\fl}{\operatorname{fl}}
\newcommand{\down}{\operatorname{down}}
\newcommand{\up}{\operatorname{up}}
\newcommand{\abs}[1]{\left|#1\right|}
\newcommand{\norminf}[1]{\left\|#1\right\|_{\infty}}
\newcommand{\GammaE}{\Gamma}
\newcommand{\eps}{\varepsilon}
\newcommand{\epsv}{\boldsymbol{\varepsilon}}
\newcommand{\qv}{\mathbf{q}}
\newcommand{\pv}{\mathbf{p}}
\newcommand{\yv}{\mathbf{y}}
\newcommand{\xv}{\mathbf{x}}
\newcommand{\one}{\mathbf{1}}
\newcommand{\zero}{\mathbf{0}}
\title{\textbf{AutoIA: Mathematical Theory and Certification Architecture}\\
A Self-Contained Specification Through Version 0.21.0.dev1}
\author{}
\date{August 2026}
\begin{document}
\maketitle

\begin{abstract}
AutoIA is a research implementation of automatic validated interval analysis for
existing engineering Python/NumPy programs.  The central objective is to accept a
deterministic engineering program, replace selected scalar and numerical operations
by validated counterparts, and return output enclosures that contain the exact real
semantics of every supported execution path for every input in a user-supplied
interval box.  The system is not based on naive line-by-line interval arithmetic
alone.  It combines explicit polynomials in the original uncertain inputs,
interval remainders, shared auxiliary error generators, semantic certification of
linear systems, path-aware execution, monotonicity and Bernstein range analysis,
outward-rounded binary64 proof kernels, and prepared repeated-query certificate
reuse.

This document is a mathematical specification of the current design through
AutoIA version 0.21.0.dev1.  Every symbol used below is defined explicitly.  The
document separates \emph{proposal calculations}, which may use ordinary floating
point and may be inaccurate, from \emph{proof calculations}, whose only purpose is
to establish inclusion.  It also states the major current limitations and the
precise sense in which prepared-session restriction can answer a new interior
input box without re-executing the engineering model.
\end{abstract}

\tableofcontents
\newpage

%===============================================================================
\section{Problem statement and certification target}
%===============================================================================

Let
\[
  \mathcal{P}:\R^m\rightarrow\R^r
\]
denote the mathematical semantics of a supported deterministic engineering
program.  The symbol \(m\) denotes the number of uncertain scalar input
parameters and \(r\) denotes the number of scalar outputs.  The input vector is
\[
  \pv=(p_1,\ldots,p_m)^T\in\R^m,
\]
and the output vector is
\[
  \yv=(y_1,\ldots,y_r)^T=\mathcal{P}(\pv)\in\R^r.
\]

The user supplies an interval for every uncertain input,
\[
  p_i\in I_i=[\underline p_i,\overline p_i],
  \qquad i=1,\ldots,m,
\]
where \(\underline p_i\le \overline p_i\) are finite binary64 numbers interpreted
as exact real interval endpoints.  The complete input box is the Cartesian
product
\[
  \Pin=I_1\times\cdots\times I_m\subset\R^m.
\]

The exact reachable output set is
\[
  \Yset^\star
  =\left\{\mathcal{P}(\pv):\pv\in\Pin\right\}.
\]
AutoIA does not promise to compute the exact interval hull of this set.  The
primary contract is inclusion.  For every scalar output \(y_j\), AutoIA returns
an interval
\[
  Y_j=[\underline y_j,\overline y_j]
\]
such that
\[
  \boxed{
  y_j(\pv)\in Y_j
  \quad\text{for every }\pv\in\Pin.
  }
\]
For a piecewise/path-aware result, the system may instead return a finite union
of certified branch enclosures.  The reported global lower and upper bounds are
then the interval hull of that finite union.

\subsection{What is meant by program semantics}

For generic supported scalar arithmetic, stored binary64 constants in the user
program are interpreted as exact real constants equal to their stored values.
Recognized mathematical primitives are interpreted according to their documented
real mathematical semantics.  Examples include exact real \(\exp\), \(\log\),
\(\sin\), \(\cos\), \(\operatorname{erf}\), normal PDF/CDF, and the exact solution
of a recognized linear system \(A x=b\).  The semantic linear-solve primitive
therefore certifies the mathematical solution of the represented matrix family,
not the floating-point rounding behavior of NumPy's implementation of
\texttt{linalg.solve}.

If the user manually codes an iterative or approximate algorithm that is not a
recognized semantic primitive, the default goal is to certify the supported
program semantics as written.  AutoIA does not silently replace an arbitrary
legacy algorithm by a different mathematical problem merely because the latter
looks like the likely engineering intent.

%===============================================================================
\section{Notation used throughout the document}
%===============================================================================

This section defines the notation that appears later.

\begin{longtable}{p{0.20\textwidth}p{0.73\textwidth}}
\toprule
Symbol & Definition \\
\midrule
\endhead
\(m\) & Number of uncertain scalar physical inputs. \\
\(r\) & Number of scalar outputs. \\
\(\pv\) & Physical input vector \((p_1,\ldots,p_m)^T\). \\
\(\Pin\) & User-supplied Cartesian input box. \\
\(p_{c,i}\) & Stored binary64 center chosen for input interval \(I_i\). \\
\(q_i\) & Centered physical coordinate \(q_i=p_i-p_{c,i}\). \\
\(h_i\) & Rigorous nonnegative half-width satisfying \(q_i\in[-h_i,h_i]\). \\
\(\qv\) & Centered physical vector \((q_1,\ldots,q_m)^T\). \\
\(\Qdom\) & Symmetric centered box \(\prod_{i=1}^m[-h_i,h_i]\). \\
\(\alpha\) & Multi-index \((\alpha_1,\ldots,\alpha_m)\) with nonnegative integer components. \\
\(|\alpha|\) & Total degree \(\alpha_1+\cdots+\alpha_m\). \\
\(\qv^\alpha\) & Monomial \(\prod_{i=1}^m q_i^{\alpha_i}\). \\
\(e_i\) & Unit multi-index whose \(i\)-th component is one and all others are zero. \\
\(P(\qv)\) & Stored point-coefficient polynomial in the original centered physical inputs. \\
\(R\) & Interval remainder, usually written \([\underline R,\overline R]\). \\
\(s\) & Number of active auxiliary generator variables. \\
\(\epsv\) & Auxiliary vector \((\eps_1,\ldots,\eps_s)^T\) with each \(\eps_j\in[-1,1]\). \\
\(G\) & Generator coefficient matrix.  For a scalar, this reduces to a row/vector of coefficients. \\
\(\Gamma_E\) & Concretization/interpretation of a validated representation \(E\). \\
\(A(\qv)\) & Parameter-dependent matrix in a semantic linear system. \\
\(b(\qv)\) & Parameter-dependent right-hand side in a semantic linear system. \\
\(X(\qv)\) & Stored polynomial predictor for a linear-system solution. \\
\(C\) & Ordinary floating-point point preconditioner used in a linear-system certificate. \\
\(B\) & Nonnegative matrix satisfying a rigorous componentwise defect bound. \\
\(c\) & Nonnegative vector satisfying a rigorous preconditioned residual bound. \\
\(\rho\) & Rigorous upper bound on \(\|B\|_\infty\). \\
\(\|B\|_\infty\) & Maximum absolute row-sum norm \(\max_i\sum_j |B_{ij}|\); since \(B\ge0\), this is \(\max_i\sum_jB_{ij}\). \\
\(\Hull(S)\) & Smallest interval containing a scalar set \(S\), or componentwise interval hull for a vector set. \\
\(\down(x)\) & A binary64 lower endpoint known to be no greater than the exact real quantity \(x\). \\
\(\up(x)\) & A binary64 upper endpoint known to be no less than the exact real quantity \(x\). \\
\bottomrule
\end{longtable}

Componentwise vector/matrix inequalities such as
\[
  |v|\le c
\]
mean \(|v_i|\le c_i\) for every component \(i\).  Similarly,
\[
  |M|\le B
\]
means \(|M_{ij}|\le B_{ij}\) for every entry.

%===============================================================================
\section{Rigorous centering of the input box}
%===============================================================================

Polynomial representations use centered coordinates.  For each user interval
\[
  I_i=[\underline p_i,\overline p_i],
\]
AutoIA chooses a finite binary64 proposal center \(p_{c,i}\).  Conceptually this
is near the arithmetic midpoint, but the proof does not require the chosen center
to be the exact midpoint.  The trusted backend then rigorously encloses
\[
  I_i-p_{c,i}
\]
using outward-rounded interval subtraction.  If the resulting interval is
\([\ell_i,u_i]\), AutoIA defines
\[
  h_i=\max(|\ell_i|,|u_i|).
\]
Therefore
\[
  I_i\subseteq[p_{c,i}-h_i,p_{c,i}+h_i],
\]
and hence every physical input in \(I_i\) admits the centered coordinate
\[
  q_i=p_i-p_{c,i}\in[-h_i,h_i].
\]
The complete centered domain is
\[
  \Qdom=\prod_{i=1}^m[-h_i,h_i].
\]

The symmetric centered domain can be slightly larger than the original user
interval because outward rounding may enlarge a boundary by one or more ulps.
This enlargement is sound but matters near function-domain boundaries.  For
example, a user interval whose exact lower endpoint is zero may produce a
centered polynomial domain whose rigorous reconstructed lower endpoint is a tiny
negative number.  AutoIA must not silently clip that negative number when
certifying \(\sqrt{x}\) or \(\log(x)\); the current implementation fails closed
in such a case unless a representation/path/domain refinement proves the required
real domain.

%===============================================================================
\section{The representation-independent enclosure invariant}
%===============================================================================

Every validated scalar value is associated with a representation object \(E\)
and a concretization map
\[
  \Gamma_E(\qv)\subseteq\R.
\]
The fundamental invariant is
\[
  \boxed{
  z(\qv)\in\Gamma_E(\qv)
  \quad\text{for all }\qv\in\Qdom,
  }
\]
where \(z\) is the exact real semantic quantity represented by \(E\).

Suppose an exact mathematical operation is denoted
\[
  z=\operatorname{op}(z_1,\ldots,z_k).
\]
A validated implementation \(\operatorname{op}^{\#}\) is sound when
\[
  \boxed{
  \Gamma_{\operatorname{op}^{\#}(E_1,\ldots,E_k)}(\qv)
  \supseteq
  \operatorname{op}\bigl(
      \Gamma_{E_1}(\qv),\ldots,\Gamma_{E_k}(\qv)
  \bigr)
  }
\]
for every \(\qv\in\Qdom\), with the right-hand side interpreted as the set of
all values of the exact operation on arguments chosen from the input sets.

Representation conversion is similarly required to be inclusion-preserving.  If
\(E_A\) is converted to \(E_B\), then
\[
  \Gamma_{E_A}(\qv)\subseteq\Gamma_{E_B}(\qv)
\]
for every \(\qv\).  Promotion priorities in the software influence strategy only;
they are not themselves part of the proof.

%===============================================================================
\section{The interval representation}
%===============================================================================

The simplest scalar representation is a single interval
\[
  E=I=[\underline z,\overline z].
\]
Its concretization is independent of \(\qv\):
\[
  \Gamma_E(\qv)=I.
\]
This representation forgets explicit dependence on the original physical inputs,
but it is robust and useful as a fallback or for calculations where interval
semantics are already natural.

The natural interval extension of an exact operation is used when supported.  For
example, if \(X=[\underline x,\overline x]\) and
\(Y=[\underline y,\overline y]\), then interval addition returns an outward
rounded enclosure of
\[
  [\underline x+\underline y,\overline x+\overline y].
\]
Interval multiplication encloses all four endpoint products.  Division requires
a denominator interval separated from zero.

%===============================================================================
\section{Point-coefficient polynomial plus interval remainder}
%===============================================================================

The principal dependency-preserving representation is
\[
  \boxed{
  z(\qv)\in P(\qv)+R,
  }
\]
where
\[
  P(\qv)=\sum_{\alpha\in\mathcal{A}} c_\alpha\qv^\alpha
\]
is a sparse polynomial with stored binary64 point coefficients \(c_\alpha\),
\(\mathcal{A}\) is a finite set of multi-indices, and
\[
  R=[\underline R,\overline R]
\]
is an interval remainder.  The concretization is
\[
  \Gamma_E(\qv)=P(\qv)+R
  =\{P(\qv)+r:r\in R\}.
\]

The stored coefficient \(c_\alpha\) is interpreted as an exact real number equal
to its binary64 value.  When a coefficient is produced by floating-point proposal
arithmetic, AutoIA separately computes a rigorous interval containing the exact
coefficient that the mathematical operation would produce, subtracts the stored
point coefficient, and ranges that coefficient error over the input domain into
\(R\).  This separation allows fast ordinary floating-point proposals while
keeping the final enclosure rigorous.

\subsection{Multi-index monomial range on a symmetric box}

For
\[
  \alpha=(\alpha_1,\ldots,\alpha_m),
\]
define
\[
  H_\alpha=\prod_{i=1}^m h_i^{\alpha_i}.
\]
The trusted implementation computes an upward-rounded upper bound for this
nonnegative product.  On the symmetric box \(\Qdom\), the monomial range is
\[
  \qv^\alpha\in
  \begin{cases}
    [0,H_\alpha], & \text{if every }\alpha_i\text{ is even},\\
    [-H_\alpha,H_\alpha], & \text{otherwise}.
  \end{cases}
\]
This parity-aware rule is sharper than blindly using \([-H_\alpha,H_\alpha]\)
for every monomial.

%===============================================================================
\section{Shared-generator polynomial representation}
%===============================================================================

The polynomial-plus-remainder representation loses correlation when the same
unstructured error is reused.  AutoIA therefore also supports
\[
  \boxed{
  z(\qv,\epsv)
  \in
  P(\qv)+\sum_{j=1}^s g_j\eps_j+R,
  \qquad \eps_j\in[-1,1].
  }
\]
For a vector of \(r\) outputs this is
\[
  \boxed{
  \mathbf{z}(\qv,\epsv)
  \in
  \mathbf{P}(\qv)+G\epsv+R,
  }
\]
where \(G\in\R^{r\times s}\) and \(R\) denotes a componentwise interval box.

The variables \(\eps_j\) are \emph{auxiliary} uncertainty variables.  They do not
represent additional physical uncertain inputs.  Their purpose is to preserve
provenance of approximation/certification error.  Two coefficients bearing the
same generator identifier refer to the same \(\eps_j\); distinct identifiers are
treated as independent coordinates in \([-1,1]^s\).

This distinction prevents an unsound shortcut.  If two unrelated calculations
happen to produce the same numerical error radius, AutoIA does not infer that
the errors are identical.  Shared provenance is created only by copying/reusing
a value or by a semantic operation that rigorously carries the same generator
through its output predictor.

\subsection{Why fan-out cancellation improves}

Suppose
\[
  y=P(\qv)+g\eps_1.
\]
If the same represented value is used twice, then
\[
  y-y
  =P-P+g\eps_1-g\eps_1
  =0.
\]
A pure interval remainder would instead behave like two independent copies of
the same interval and could fail to cancel.

\subsection{Closure under multiplication}

The representation is affine in \(\epsv\), so it is not algebraically closed
under multiplication.  If
\[
  x=P_x+G_x\epsv+R_x,
  \qquad
  y=P_y+G_y\epsv+R_y,
\]
then products such as
\[
  q_i\eps_j
  \quad\text{and}\quad
  \eps_i\eps_j
\]
appear.  The current representation keeps the constant-in-\(\qv\), first-order
in-\(\epsv\) contribution explicitly and rigorously ranges terms outside the
representation family into fresh bounded error.  Soundness is therefore retained
without pretending that \(P(\qv)+G\epsv+R\) is closed under arbitrary nonlinear
operations.

%===============================================================================
\section{Rigorous polynomial algebra}
%===============================================================================

\subsection{Addition and subtraction}

Let
\[
  x\in P_x+R_x,
  \qquad
  y\in P_y+R_y.
\]
The exact algebraic sum is
\[
  x+y\in(P_x+P_y)+(R_x+R_y).
\]
AutoIA computes point coefficient proposals for \(P_x+P_y\), rigorously encloses
coefficient-rounding errors, and adds their ranged effect to the remainder.
Subtraction is addition of the negated representation.

\subsection{Truncated multiplication}

Let the retained maximum polynomial order be \(K\).  The exact polynomial
product is
\[
  P_xP_y
  =\sum_\gamma C_\gamma\qv^\gamma,
  \qquad
  C_\gamma=\sum_{\alpha+\beta=\gamma}c_\alpha d_\beta.
\]
For every \(\gamma\) satisfying \(|\gamma|\le K\), AutoIA stores a point proposal
\(\widehat C_\gamma\).  It rigorously encloses the exact coefficient
\(C_\gamma\) and returns the difference
\[
  C_\gamma-\widehat C_\gamma
\]
as a coefficient-error interval.  That error is multiplied by a rigorous
monomial range and accumulated into the remainder.

For discarded terms with \(|\gamma|>K\), the full exact coefficient interval is
ranged over \(\Qdom\) and added to the remainder.  The user-selected order \(K\)
therefore changes sharpness and cost but not the inclusion contract.

The complete product also contains interactions with existing remainders.  These
are enclosed by ordinary interval multiplication using rigorous ranges of the
retained polynomial parts.

%===============================================================================
\section{Powers, reciprocal, and division}
%===============================================================================

\subsection{Integer powers}

For \(n\in\Z\), positive powers are constructed by validated multiplication.
For \(n<0\),
\[
  x^n=\frac{1}{x^{-n}},
\]
so AutoIA must prove that zero is excluded from the denominator range.
Negative bases are valid for integer exponents.

\subsection{Non-integer constant real powers}

For a finite constant binary64 exponent \(p\in\R\setminus\Z\), the current real
branch requires the certified base range to be strictly positive.  AutoIA may
decompose
\[
  p=n+r,
\]
where \(n\in\Z\) and \(|r|\le1/2\), then use integer powers together with
\(\sqrt{\cdot}\) for half powers or the identity
\[
  x^r=\exp(r\log x)
\]
for the remaining positive-real case.  An uncertain exponent is outside the
current supported contract.

\subsection{Reciprocal certificate}

Suppose a represented denominator \(y\) has a rigorous range
\[
  Y=[\underline Y,\overline Y]
\]
with
\[
  0\notin Y.
\]
Let \(Q\) be any stored polynomial proposal for \(1/y\).  The exact identity
\[
  \frac{1}{y}-Q
  =\frac{1-yQ}{y}
\]
allows AutoIA to bound the residual \(1-yQ\), divide by the denominator interval
\(Y\), and obtain a certified reciprocal.  The quality of \(Q\) affects only the
remainder width.

%===============================================================================
\section{Trusted binary64 interval kernel and outward rounding}
%===============================================================================

Machine rigor requires every proof-critical floating-point operation to enclose
its exact real counterpart.  AutoIA therefore distinguishes two classes of
calculation.

\begin{enumerate}[leftmargin=*]
\item \textbf{Proposal calculations.}  These may use ordinary round-to-nearest
binary64, NumPy linear algebra, heuristic scores, or approximate inverses.
Their outputs are never trusted merely because they were computed.
\item \textbf{Proof calculations.}  These establish interval inclusion and are
performed by the audited trusted backend or by operations reducible to trusted
interval primitives.
\end{enumerate}

For a scalar operation \(z=x\circ y\), the trusted backend returns binary64
endpoints
\[
  [\underline z,\overline z]
\]
known to contain the exact real value.  On the reference backend this is done by
round-to-nearest arithmetic followed by \texttt{nextafter} toward
\(-\infty\) for the lower endpoint and \(+\infty\) for the upper endpoint.

The backend construction audits the binary64 environment.  The audit checks,
among other things, the binary64 format, NumPy \texttt{float64} size,
round-to-nearest behavior, adjacent \texttt{nextafter} values, gradual
underflow/subnormals, and required vectorized arithmetic behavior.  If the
environment does not satisfy the assumptions, AutoIA fails closed.

\subsection{Why a single outward adjustment after a long reduction is insufficient}

A long dot product
\[
  s=\sum_{i=1}^n a_i b_i
\]
contains rounding error in every multiplication and every addition.  Merely
computing an ordinary BLAS result and moving it outward by one ulp does not in
general enclose the exact real dot product.

The Step-9/9b trusted fast path instead encloses every product and combines those
product intervals through a pairwise interval-addition tree.  If each leaf
contains its exact product and each internal node contains the exact sum of its
two child sets, induction proves that the root interval contains the exact dot
product.  Vectorization changes execution speed, not the mathematical proof.

The same idea is used for point-by-interval and interval-by-interval reductions,
and many independent dot-product trees may be batched simultaneously across
matrix rows/columns.

%===============================================================================
\section{Polynomial range analysis}
%===============================================================================

A validated symbolic representation is useful only if it can be converted to a
scalar interval bound.  AutoIA uses several sound range engines.

\subsection{Termwise range}

For
\[
  P(\qv)=\sum_\alpha c_\alpha\qv^\alpha,
\]
each monomial is enclosed separately and the resulting intervals are summed.
This method is inexpensive and robust, but it may lose cancellation between
terms.

\subsection{Derivative-based monotonicity}

For coordinate \(q_i\), AutoIA constructs the polynomial derivative
\[
  \frac{\partial P}{\partial q_i}
  =\sum_{\alpha:\alpha_i>0}
  \alpha_i c_\alpha\qv^{\alpha-e_i}.
\]
If a rigorous range proves
\[
  \frac{\partial P}{\partial q_i}\ge0
\]
throughout the box, then \(P\) is nondecreasing in \(q_i\); extrema occur at the
appropriate endpoint of that coordinate.  If the derivative is rigorously
nonpositive, the opposite endpoint is used.  Monotone directions can therefore
be fixed before ranging the unresolved directions.

\subsection{Bernstein enclosure}

After mapping the active box to a unit box, the polynomial can be expressed in a
multivariate Bernstein basis.  The range of a polynomial over the unit box is
contained in the interval hull of its Bernstein coefficients.  AutoIA computes
an outward-rounded version of this transformation.  Because the number of
Bernstein coefficients grows as a product of per-variable degrees, the method is
used selectively.

\subsection{Monotonicity-reduced Bernstein}

If some coordinates are already proved monotone, those coordinates are fixed at
appropriate endpoints and Bernstein conversion is applied only to unresolved
directions.  This can greatly reduce coefficient count.

\subsection{Portfolio principle}

If two independent rigorous methods produce intervals \(I_1\) and \(I_2\) that
both contain the same quantity, their intersection also contains that quantity,
provided the intersection is nonempty.  AutoIA therefore uses range-method
portfolios to improve sharpness without making correctness depend on which
method happens to be tighter.

%===============================================================================
\section{Validated elementary functions by Taylor composition}
%===============================================================================

For a scalar represented argument \(x\), AutoIA first obtains a rigorous argument
interval
\[
  X=[\underline x,\overline x].
\]
Choose a stored binary64 expansion center \(c\in X\) and write
\[
  \delta=x-c.
\]
For a smooth function \(f\), a degree-\(K\) Taylor representation is
\[
  f(x)
  =\sum_{n=0}^{K}a_n\delta^n+R_{K+1},
  \qquad
  a_n=\frac{f^{(n)}(c)}{n!}.
\]
The coefficient \(a_n\) is itself represented by a rigorous interval.  AutoIA
chooses a point proposal \(\widehat a_n\), stores it in the polynomial, and
propagates the difference between the rigorous coefficient interval and
\(\widehat a_n\) as certified error.  A rigorous analytic bound on the omitted
Taylor tail \(R_{K+1}\) is then added.

This generic composition mechanism is used for \(\exp\), \(\log\),
\(\sqrt{\cdot}\), \(\sin\), \(\cos\), \(\operatorname{erf}\), the standard normal
PDF, and the standard normal CDF, with function-specific coefficient recurrences
and tail bounds.

\subsection{Exponential}

For sufficiently small \(|r|\),
\[
  e^r=\sum_{k=0}^{N}\frac{r^k}{k!}+R_N.
\]
If \(|r|\le a<1\), one convenient rigorous tail bound is
\[
  |R_N|
  \le
  \frac{a^{N+1}}{(N+1)!(1-a)}.
\]
The trusted scalar kernel repeatedly halves a rigorous argument interval until it
is small enough for the series, then reconstructs the result by trusted squaring.

\subsection{Logarithm}

For \(x>0\), binary exponent extraction produces
\[
  x=2^k m,
\]
with a scaled mantissa \(m\).  Then
\[
  \log x=k\log2+\log m.
\]
Using
\[
  z=\frac{m-1}{m+1},
\]
we have
\[
  \log m=2\operatorname{atanh}(z),
\]
which is evaluated by a rigorously bounded odd series.  The constant \(\log2\)
is itself generated rigorously rather than trusted from a platform library.

\subsection{Sine and cosine}

A rigorous interval for \(\pi\) is generated from Machin's formula.  An ordinary
floating-point integer reduction index may be proposed, but the proof forms the
reduced argument with the rigorous \(\pi\) interval.  Taylor sine/cosine bounds
are used when the reduced interval is small enough; otherwise a safe range such
as \([-1,1]\) is available.

\subsection{Error function and normal distribution functions}

The error function is
\[
  \operatorname{erf}(x)
  =\frac{2}{\sqrt\pi}\int_0^x e^{-t^2}\,dt.
\]
Moderate positive arguments use a bounded alternating Maclaurin series.  Larger
arguments use the complementary error function and an integration-by-parts
asymptotic bracket.  Odd symmetry handles negative arguments.

The standard normal density and CDF are
\[
  \phi(x)=\frac{1}{\sqrt{2\pi}}e^{-x^2/2},
\]
\[
  \Phi(x)=\frac12\left(1+\operatorname{erf}\left(\frac{x}{\sqrt2}\right)\right).
\]
Their polynomial Taylor coefficients are generated using differential recurrences,
including
\[
  \phi'(x)+x\phi(x)=0
\]
and
\[
  \Phi'(x)=\phi(x).
\]
A normal distribution with constant mean \(\mu\) and standard deviation
\(\sigma>0\) is handled using
\[
  z=\frac{x-\mu}{\sigma},
  \qquad
  f_X(x)=\frac{\phi(z)}{\sigma},
  \qquad
  F_X(x)=\Phi(z).
\]

%===============================================================================
\section{Semantic certification of parameter-dependent linear systems}
%===============================================================================

Consider a represented family of linear systems
\[
  A(\qv)x(\qv)=b(\qv),
\]
where \(A(\qv)\in\R^{n\times n}\), \(x(\qv)\in\R^n\), and
\(b(\qv)\in\R^n\).  The matrix and right-hand side may have retained polynomial
coefficients of arbitrary supported order plus rigorous remainders.

\subsection{General polynomial solution predictor}

Write
\[
  A(\qv)\approx\sum_\beta A_\beta\qv^\beta,
  \qquad
  b(\qv)\approx\sum_\beta b_\beta\qv^\beta,
\]
and propose
\[
  X(\qv)=\sum_{|\alpha|\le K}x_\alpha\qv^\alpha.
\]
Formal coefficient matching gives the recurrence
\[
  \boxed{
  A_0x_\alpha
  =
  b_\alpha
  -
  \sum_{\substack{\beta\ne0\\\beta\le\alpha}}
      A_\beta x_{\alpha-\beta}.
  }
\]
Here \(\beta\le\alpha\) means \(\beta_i\le\alpha_i\) for every coordinate.
The recurrence is only a predictor generator.  AutoIA may use an ordinary
floating-point solve, a different reference matrix, a family-hull midpoint, or
another finite proposal.

\subsection{Residual identity for an arbitrary predictor}

Let \(X\) be whatever predictor is actually stored.  Define the true solution
error
\[
  e=x-X.
\]
Define the residual
\[
  r=b-AX.
\]
Because \(Ax=b\),
\[
  A(X+e)=b,
\]
so
\[
  \boxed{Ae=r.}
\]
This identity is exact and does not depend on how \(X\) was obtained.

Choose any finite point matrix \(C\) as a preconditioner.  Then
\[
  CAe=Cr.
\]
Add and subtract \(e\):
\[
  e-(I-CA)e=Cr,
\]
and hence
\[
  \boxed{
  e=Cr+(I-CA)e.
  }
\]

Suppose trusted calculations establish a nonnegative vector \(c\) and
nonnegative matrix \(B\) satisfying, for every admissible input,
\[
  |Cr|\le c,
\]
\[
  |I-CA|\le B.
\]
Taking componentwise absolute values gives
\[
  |e|\le c+B|e|.
\]

\subsection{Contraction theorem}

Define
\[
  \rho=\|B\|_\infty=\max_i\sum_jB_{ij}.
\]
If a rigorous upper bound proves
\[
  \boxed{\rho<1,}
\]
then the family is nonsingular throughout the represented domain and
\[
  \boxed{
  |e|\le(I-B)^{-1}c.
  }
\]
A simpler scalar consequence is obtained by defining
\[
  \eta=\|c\|_\infty,
  \qquad
  \delta=1-\rho>0.
\]
Then
\[
  \|e\|_\infty\le\frac{\eta}{\delta}.
\]
The implementation computes an outward-rounded upper bound for this ratio and
may subsequently refine the componentwise correction by the monotone iteration
\[
  d^{(k+1)}=\min\left(d^{(k)},\;\up(c+Bd^{(k)})\right),
\]
starting from a safe scalar vector bound.

\subsection{Preconditioner portfolio}

The center matrix inverse is a useful proposal but can be a poor contraction
preconditioner.  AutoIA also constructs a rigorous interval hull of the matrix
family, chooses a finite representative such as its componentwise midpoint, and
proposes the inverse of that representative.  Every candidate is checked against
the complete matrix family.  Therefore a midpoint that is not physically
realizable still yields a sound certificate if the final defect bound passes.

%===============================================================================
\section{Generator-aware semantic linear solve}
%===============================================================================

Suppose the solve inputs contain shared auxiliary generators:
\[
  A(\qv,\epsv)
  \in
  \widehat A(\qv)+\sum_{j=1}^s A_j\eps_j+R_A,
\]
\[
  b(\qv,\epsv)
  \in
  \widehat b(\qv)+\sum_{j=1}^s b_j\eps_j+R_b.
\]
AutoIA proposes a solution of the form
\[
  \boxed{
  X(\qv,\epsv)
  =P(\qv)+\sum_{j=1}^s u_j\eps_j.
  }
\]
For generator \(\eps_j\), a useful sensitivity proposal is obtained from a
reference matrix \(A_{\mathrm{ref}}\) by
\[
  A_{\mathrm{ref}}u_j
  =b_j-A_jx_0,
\]
where \(x_0=P(0)\).  Again, this is only a proposal.

The proof then evaluates the actual residual
\[
  r(\qv,\epsv)
  =b(\qv,\epsv)-A(\qv,\epsv)X(\qv,\epsv)
\]
over the joint domain
\[
  \qv\in\Qdom,
  \qquad
  \epsv\in[-1,1]^s,
\]
and uses the same fixed-point theorem.  Incoming generator identifiers are
preserved in the explicit predictor when justified.  The remaining componentwise
correction proven by the contraction theorem is assigned fresh independent
auxiliary generators because the theorem proves only a box for the residual
error, not a stronger coupled geometry.

%===============================================================================
\section{Path-aware execution and uncertain control flow}
%===============================================================================

Consider a scalar predicate such as
\[
  \texttt{lhs < rhs}.
\]
AutoIA forms the validated difference
\[
  d=\texttt{lhs}-\texttt{rhs}
\]
and computes a rigorous range
\[
  d\in[\underline d,\overline d].
\]
For the strict less-than relation:
\begin{itemize}[leftmargin=*]
\item if \(\overline d<0\), the predicate is proved true throughout the current
box;
\item if \(\underline d\ge0\), it is proved false throughout the current box;
\item otherwise the predicate is unresolved.
\end{itemize}
Analogous exact-real rules are used for \(\le,>,\ge,==,!=\) and scalar truth
checking.

If the predicate is unresolved, AutoIA may split the physical input box and
re-execute.  For an exact affine one-variable predicate
\[
  d(p_i)=a+bp_i,
  \qquad b\ne0,
\]
the exact root
\[
  p_i^\star=-a/b
\]
can be used as a direct partition point when it is rigorously recognized.

For a general nonlinear or multivariate decision surface, finitely many Cartesian
boxes may not align with the exact path boundary.  If subdivision cannot resolve
the predicate under the allowed budget, AutoIA conservatively enumerates both
feasible Boolean outcomes and certifies the continuation of each path separately.
If the two branch enclosures are \(Y_T\) and \(Y_F\), then
\[
  \mathcal{P}(\Pin)\subseteq Y_T\cup Y_F.
\]
The branch formula may be evaluated on a superset of its true feasible domain;
this can make the result wider or cause a boundary singularity to fail
certification, but it cannot omit a true output.

Path-local continuation is essential.  AutoIA does not immediately hull two
branch values and then continue the rest of the program.  Instead every feasible
path carries its own validated values through downstream arithmetic, elementary
functions, and semantic solves.  The union is formed only at the final piecewise
result.

%===============================================================================
\section{Adaptive subdivision controller}
%===============================================================================

Subdivision is a sharpness/resource strategy, not a correctness prerequisite.
Every retained child box is independently certified.  If a controller chooses a
poor direction, the result may be slower or wider but does not become unsound.

The practical split score introduced in version 0.19.0.dev1 uses the already
certified parent polynomial.  For candidate coordinate \(q_d\), a half-box child
has a shifted center
\[
  s_d=\pm h_d/2
\]
and new coordinate \(r_d\) satisfying
\[
  q_d=r_d+s_d,
  \qquad
  r_d\in[-h_d/2,h_d/2].
\]
The binomial identity
\[
  q_d^{\alpha_d}
  =\sum_{b=0}^{\alpha_d}
    {\alpha_d\choose b}s_d^{\alpha_d-b}r_d^b
\]
permits a cheap predicted re-centering without executing the engineering model.
The score estimates derivative-sign resolution, reduced-Bernstein burden, and
output-width gain.  Only the selected child pair is actually certified.  The
controller then checks realized structural/width improvement before accepting
the split.

The old expensive one-step lookahead remains available as a research oracle.
Because controller decisions are outside the trusted proof kernel, agreement with
that oracle is an efficiency/sharpness validation issue rather than a soundness
condition.

%===============================================================================
\section{Prepared sessions and repeated-query certification}
%===============================================================================

Version 0.21.0.dev1 separates static program preparation from repeated box
certification.  Let one fixed engineering program \(\mathcal{P}\) and one fixed
AutoIA configuration be used for a sequence of boxes
\[
  \Pin^{(1)},\Pin^{(2)},\ldots.
\]
A prepared session reuses the local semantic-dispatch clone and one audited
trusted backend.  More importantly, it stores previously certified boxes and
results in a bounded certificate cache.

\subsection{Exact repeated-box reuse}

If a new query box has exactly the same stored binary64 endpoints as a cached
box, the previous certificate already proves the same statement.  No engineering
model execution is required.

\subsection{Restriction of a parent certificate to an interior subbox}

Suppose a parent box \(D\) has already been certified and a new query box \(Q\)
satisfies
\[
  Q\subseteq D.
\]
Set-theoretically, every certificate valid on \(D\) is automatically valid on
\(Q\).  For a polynomial certificate, AutoIA can do better than returning the
same parent output interval: it re-centers and re-ranges the stored polynomial on
\(Q\).

Let the old centered coordinate be
\[
  \qv=\pv-\mathbf{c}_{\mathrm{old}},
\]
and the new centered coordinate be
\[
  \mathbf{r}=\pv-\mathbf{c}_{\mathrm{new}}.
\]
Then exactly
\[
  \qv=\mathbf{r}+\mathbf{s},
\]
where
\[
  \mathbf{s}=\mathbf{c}_{\mathrm{new}}-\mathbf{c}_{\mathrm{old}}.
\]
For a monomial
\[
  \qv^\alpha=\prod_{i=1}^m(r_i+s_i)^{\alpha_i}.
\]
Using the multivariate binomial expansion,
\[
  \qv^\alpha
  =
  \sum_{\beta\le\alpha}
  \left[
  \prod_{i=1}^m
  \binom{\alpha_i}{\beta_i}
  s_i^{\alpha_i-\beta_i}
  \right]
  \mathbf{r}^\beta.
\]
Therefore if
\[
  P(\qv)=\sum_\alpha c_\alpha\qv^\alpha,
\]
the exact re-centered coefficient of \(\mathbf{r}^\beta\) is
\[
  c'_\beta
  =
  \sum_{\alpha\ge\beta}
  c_\alpha
  \prod_{i=1}^m
  \binom{\alpha_i}{\beta_i}
  s_i^{\alpha_i-\beta_i}.
\]

The exact shift \(s_i\) is not assumed to equal the ordinary floating-point
subtraction of the two stored centers.  The trusted backend forms an interval
\(S_i\) containing the exact real difference.  Every transformed coefficient is
computed as a rigorous interval using \(S_i\), while an ordinary binary64 point
coefficient is stored as a proposal.  The difference between the exact coefficient
interval and the stored proposal is rigorously ranged over the new centered box
and accumulated into an extra remainder.

Hence the re-centered value remains a valid certificate on the new domain.
Shared-generator coefficients are copied unchanged because the coordinate shift
acts on the physical polynomial \(P\), not on the auxiliary variables \(\epsv\).

\subsection{Containment check for the actual centered query domain}

The user query \(Q\) is not quite the domain used by the polynomial executor.
The executor constructs an outward symmetric centered envelope around each new
binary64 center.  Let that actual represented domain be \(\widetilde Q\).
Certificate restriction is used only when AutoIA proves
\[
  \widetilde Q\subseteq D.
\]
This prevents the known one-ulp symmetric-centering enlargement from causing a
parent certificate to be applied outside the domain on which it was proved.  If
this containment check fails, AutoIA performs a fresh certification instead.

\subsection{Many-box execution}

For a collection of \(N\) boxes, prepared execution currently supports serial,
thread-chunked, and process-chunked modes.  Process mode creates one prepared
session and one independently audited trusted backend in every worker process.
A streaming interface keeps only a bounded number of chunks in flight, so the
user does not need to store all \(N\) boxes or all \(N\) full result objects in
memory simultaneously.

This is an execution architecture for large repeated-query workloads.  It is not
yet a single dense numeric kernel with a leading \(N\)-box array dimension.

%===============================================================================
\section{Program-level soundness theorem}
%===============================================================================

We can summarize the architecture as a compositional theorem.

\paragraph{Assumptions.}
Assume:
\begin{enumerate}[leftmargin=*]
\item the trusted binary64 backend's audited machine assumptions hold;
\item every proof-critical interval operation returns an enclosure of its exact
real counterpart;
\item every supported validated scalar/semantic operation satisfies its enclosure
contract;
\item every representation conversion is inclusion-preserving;
\item every retained path/subdivision child is certified over a domain whose
union covers the relevant user input domain;
\item every semantic linear solve accepted by the system satisfies its rigorous
contraction condition; and
\item the program remains within the documented supported-language semantics.
\end{enumerate}

\paragraph{Conclusion.}
For every user input \(\pv\in\Pin\), the exact real output of the supported
program semantics belongs to the finite union of the validated output sets
returned by AutoIA.  Consequently the reported global lower/upper arrays contain
every exact scalar output.

\paragraph{Proof sketch.}
The proof is by induction through execution.  Initial uncertain inputs are
represented by enclosures containing every physical value in the input box.
Every supported operation maps valid input enclosures to an output enclosure
containing the exact operation on all admissible inputs.  Semantic solves are
validated independently by their residual/contraction theorem.  If execution
branches, the path engine retains every feasible outcome or subdivides the
physical domain into certified children whose union covers the original domain.
Thus the invariant holds on every active path.  Final range methods enclose each
representation's concretization, and the finite union/hull of path results
contains the exact output set.

Prepared-session restriction does not alter the theorem: a parent certificate is
reused only on a subdomain proved to lie inside the parent's certified domain,
and the coordinate transform itself is rigorously enclosed.

%===============================================================================
\section{What ordinary floating-point calculations are allowed to do}
%===============================================================================

A recurring AutoIA design principle is that an approximation may be poor without
making the final result unsound.  Ordinary floating point is therefore acceptable
for:
\begin{itemize}[leftmargin=*]
\item polynomial coefficient proposals;
\item nominal/reference linear solves;
\item approximate inverses and preconditioners;
\item midpoint representatives of interval matrix families;
\item branch split scores and controller heuristics;
\item candidate Taylor centers;
\item integer argument-reduction proposals for trigonometric functions;
\item cache-choice heuristics such as choosing the smallest containing parent
certificate.
\end{itemize}

Such calculations may affect sharpness, performance, or whether a proof attempt
succeeds.  They must not be used as the final reason an inclusion is accepted.
The proof always returns to rigorous residual, interval, range, or domain checks.

%===============================================================================
\section{Important limitations and non-guarantees in version 0.21.0.dev1}
%===============================================================================

The current system is a substantial research prototype, not a proof assistant for
arbitrary Python.  Important limitations include the following.

\begin{enumerate}[leftmargin=*]
\item \textbf{Exact hulls are not promised.}  Returned intervals can be wider
than the exact range.
\item \textbf{Supported-language boundary.}  Arbitrary Python side effects,
uncertain array indices, uncertain Boolean masks, general \texttt{np.where}, and
general uncertainty-dependent loop termination are not part of the frozen
contract.
\item \textbf{Boundary-domain over-enclosure.}  Symmetric polynomial centering
can enlarge an exact physical boundary slightly.  This can make a mathematically
valid \(\sqrt{\cdot}\) or \(\log\) query fail closed near zero.
\item \textbf{Shared-generator family is affine in auxiliary variables.}  Terms
such as \(q_i\eps_j\) and \(\eps_i\eps_j\) are bounded back into fresh error rather
than retained explicitly as a richer structure.
\item \textbf{Generator compression is not yet implemented.}  The number of
auxiliary IDs can grow in long calculations.
\item \textbf{Piecewise certificate restriction is limited.}  Exact cache reuse
works for pathwise results, but sharp geometric restriction of arbitrary cached
piecewise branch domains is not yet a general many-box feature.
\item \textbf{Large many-box execution is not yet a true vectorized polynomial
batch kernel.}  Current process parallelism and certificate reuse are useful, but
unrelated boxes are still separate certifications.
\item \textbf{GPU proof execution is not implemented.}  The public prepared/batch
architecture is designed to permit future CPU/GPU numeric-buffer backends.
\item \textbf{External semantic primitives are limited.}  Linear solve is the
major structured semantic operation.  General nonlinear roots, ODE/DAE solvers,
eigensystems, optimizers, and arbitrary legacy algorithm recognition are future
work.
\end{enumerate}

%===============================================================================
\section{A compact end-to-end mathematical example}
%===============================================================================

Consider a two-parameter model
\[
  p_1\in[0.9,1.1],
  \qquad
  p_2\in[-0.1,0.1],
\]
with
\[
  A(\pv)=
  \begin{bmatrix}
  2+p_1^2 & 0.2\\
  -0.1p_2 & 1.5+0.05p_1
  \end{bmatrix},
\]
\[
  b(\pv)=
  \begin{bmatrix}
  1+0.1p_2\\
  0.5+0.05p_1^2
  \end{bmatrix},
\]
and let \(u(\pv)\) solve
\[
  A(\pv)u(\pv)=b(\pv).
\]
Suppose the program returns
\[
  y_1=u_1,
  \qquad
  y_2=e^{u_2},
  \qquad
  y_3=u_1-u_1.
\]

AutoIA first constructs rigorous centered coordinates for \(p_1,p_2\).  Polynomial
arithmetic retains the quadratic occurrences of \(p_1\) up to the chosen order.
The semantic solve constructs a polynomial (and, in generator mode, auxiliary
sensitivity) predictor, then proves a contraction bound on the complete matrix
family.  The exponential is composed by the validated Taylor machinery.  Shared
generator provenance can preserve exact downstream cancellation in \(y_3\), so
\(y_3\) can certify as zero even when intermediate approximations carry error.

If a prepared session first certifies the master box and later receives a smaller
interior box, the stored output polynomial certificates may be re-centered using
the exact binomial coordinate transform with rigorous coefficient-error capture.
The model and the semantic solve need not run again if the new centered domain is
proved to lie inside the parent domain.

%===============================================================================
\section{Summary of the current mathematical architecture}
%===============================================================================

The current architecture can be summarized by the chain
\[
\boxed{
\begin{array}{c}
\text{user interval box}\\
\downarrow\\
\text{rigorous centered physical coordinates }\qv\\
\downarrow\\
P(\qv)+R\quad\text{or}\quad P(\qv)+G\epsv+R\\
\downarrow\\
\text{validated generic operations + semantic primitives}\\
\downarrow\\
\text{path-aware piecewise execution when required}\\
\downarrow\\
\text{termwise / monotonicity / Bernstein range certification}\\
\downarrow\\
\text{rigorous binary64 output intervals}
\end{array}}
\]
with a prepared-session layer that can reuse the static dispatch architecture and,
more importantly, restrict already-proved certificates to eligible interior
query boxes.

The central theoretical thesis is that sharp automatic interval analysis for
engineering code is best treated as a structured validated-execution problem:
retain original-input dependence explicitly, introduce auxiliary variables only
when their provenance is proved, treat important numerical operations according
to their mathematical semantics, and confine correctness-critical floating-point
work to a small audited enclosure kernel.  Performance strategies are then free
to be aggressive because any proposed approximation must still pass an independent
rigorous certificate before it is accepted.

\end{document}
